\begin{samepage}
\section{模型假设}

根据题给条件，作如下假设：

\begin{enumerate}
  \item \textbf{介质与几何。}药材为均匀、各向同性的连续介质，初始温湿场均匀且轴对称；主模型仅考虑中截面径向传递，端面影响另用二维模型检验。
  \item \textbf{传递机制。}热量与水分分别由导热和 Fick 扩散传递，耦合仅通过题给状态相关物性体现；忽略内热源、辐射、内部宏观对流、重力渗流、托盘接触传热及汽化潜热。
  \item \textbf{表面交换。}表面采用对流换热、传质边界，四问均取 $h_T=25~\mathrm{W/(m^2K)}$、$h_m=8\times10^{-7}~\mathrm{m/s}$。
  \item \textbf{环境延拓。}超过附件一的 14400~s 后，环境温度和水分浓度保持末值 $50.165~{}^\circ\mathrm C$、$0.04986~\mathrm{kg/kg}$。
  \item \textbf{收缩闭合。}问题四长度不变，材料点沿径向同比例收缩；干骨架质量守恒且密度在截面内均匀，满足 $\rho_dR^2=\rho_{d0}R_0^2$。题给 $\rho c_p$ 仅作有效储热系数。
  \item \textbf{浓度口径。}药材与环境水分浓度均按同一干基比值处理，直接用于扩散方程及表面传质边界。
\end{enumerate}
\end{samepage}
